
\sourcenote{Genre: admissible rule (tier 2) --- a proof-tree-rewriting
recipe. The tier taxonomy defines tier 2 as \emph{verified unpacking
into core rules} (one node $\to$ many), with rules admissible by
\emph{reduction to verified machinery} as the recorded boundary case;
\rulename{Norm} is of the latter kind, and this section --- the book's
first admissible rule --- states the general template that the later
recipes instantiate.}

\begin{brecipe}[the admissible template, general form]\label{rec:norm-template}
An admissible rule $R$ with premises $P_1, \dots, P_n$ and conclusion $C$
is warranted by a \emph{rewrite}: a function taking any proof-tree node
\[
\inferrule*[Right=$R$]{\;\mathcal{T}_1 \rhd P_1 \\ \cdots \\ \mathcal{T}_n \rhd P_n\;}{C}
\qquad\rightsquigarrow\qquad
\mathcal{D}\bigl[\mathcal{T}_1, \dots, \mathcal{T}_n\bigr] \rhd C
\]
where $\mathcal{D}$ is a derivation fragment over core (and
previously-justified admissible) rules with the same conclusion,
consuming the premise subtrees $\mathcal{T}_i$ at its leaves. Applying
the rewrite leaf-first eliminates every $R$-node from a finite tree in
finitely many steps, so any judgement derivable with $R$ is derivable
without it --- admissibility --- and the calculus's soundness theorem is
conserved with no new kernel semantic content. The verified object
shipped per admissible rule is either $\mathcal{D}$ itself or
(equivalently, and how this catalogue's Lean side does it) a soundness
lemma for $R$ that factors through already-verified machinery.
\end{brecipe}

\begin{brecipe}[the \rulename{Norm} instance --- the degenerate rewrite]
\label{rec:norm-instance}
\rulename{Norm} has no proof premises ($n = 0$): the node
\[
\inferrule*[Right=\rulename{Norm}]{
  \mathrm{polynomialCanonicalEq}(t_1, t_2) \\
  \text{\upshape both sides shape-stable and finitely typed}
}{\Phi \vdash t_1 = t_2}
\qquad\rightsquigarrow\qquad
\text{\upshape the checker replays the canonical forms;}\;
\text{\texttt{sound\_norm} stands as warrant.}
\]
The side condition is \emph{computational}: both sides idx-free, and
their canonical monomial polynomials \emph{identical} over $\Q$ ---
coefficient-for-coefficient, decidable, replayed by the checker. The
catalogued stratum records an unpacking sketch through
\rulename{BernCone} and \rulename{ArithAtom} (an equality as a pair of
zero-polynomial bounds); the shipped warrant is direct and stronger ---
one verified semantic lemma at all tensor shapes --- with the payload's
well-typedness fields carrying exactly what the sketch left implicit.
\end{brecipe}

\begin{bprop}[semantic warrant]\label{prop:norm-sound}
Suppose $t_1, t_2$ are idx-free with $t_1.\mathrm{toPoly} =
t_2.\mathrm{toPoly}$, the two sides carry the same defined broadcast
shape in $H$, and every wire of either side is finitely typed. Then
$\Phi \vdash t_1 = t_2$ is semantically valid over $H$ for
\emph{arbitrary} $\Phi$.
\end{bprop}
\begin{proof}
Fix any admissible valuation $\sigma$; the premise is not consulted ---
this is \emph{the \rulename{Norm} convention}, shared by every
premise-free rule of the catalogue: a fact holding at every valuation
holds under any premise. The shape and finiteness fields give both
sides defined evaluations, as tensors of the common broadcast shape:
finitely typed wires always read, and the shared shape makes the
broadcast joins on both sides agree. At each entry, evaluation is a
ring homomorphism from canonical polynomials into $\Q$ (at the
entry-wise read-out of the wires); a syntactic identity of canonical
forms is an identity of polynomials, and ring homomorphisms preserve
polynomial identities --- so the two evaluated tensors agree entry by
entry, and their shapes agree by hypothesis. Hence the equality atom
holds at $\sigma$. That is the whole proof; \texttt{sound\_norm}
packages it once, per-entry, at every tensor shape.
\end{proof}

\begin{bremark}[what broke without the typing fields]
\label{rem:norm-b13}
As first staged --- canonical-form equality alone --- the rule was
\emph{false}, twice over, and the counterexamples are why the fields
exist. (i) At a wire of extended dtype holding an infinite value, even
$\vdash t = t$ is unsatisfiable: evaluation is over exact rationals,
an infinite entry refuses, and an atom that never evaluates never
holds; since $\Phi$ is arbitrary, the staged rule derived an
unsatisfiable conclusion from a satisfiable premise. (ii) A cancelling
sub-term leaves the canonical form unchanged but not the evaluation:
$t_2 = t_1 + (w_2 - w_2)$ has $t_2.\mathrm{toPoly} =
t_1.\mathrm{toPoly}$, yet if $w_2$'s shape does not join with $t_1$'s,
the right-hand atom side never evaluates. The shared-shape and
finiteness fields exclude both; a hypothesis whose necessity is
semantic (evaluability), not algebraic, is exactly the kind a recipe
must state, or the next porter re-learns it the hard way.
\end{bremark}

\begin{bremark}[exactness discipline, and the idx-freedom gate]
\label{rem:norm-f7}
Canonical-form equality is never up-to-rescaling: $2 \cdot w$ and
$w + w$ are equal as canonical forms; $2 \cdot w$ and $w$ are not,
proportional though they are --- the matching discipline every
combination rule of this catalogue (\rulename{Farkas},
\rulename{Syzygy}) inherits when it compares atom polynomials. The
idx-freedom gate is the meaning-preservation condition of the
translation: the wire-keyed polynomial space cannot express per-entry
projections, so a projection-bearing term has no canonical form here
and the comparison is refused rather than performed on the wrong
object. (The combination rules' widened replays translate projections
into their own variable spaces; those translations belong to those
rules, not to \rulename{Norm}.)
\end{bremark}

\begin{anchors}
Recipe \ref{rec:norm-instance} warrant & \texttt{sound\_norm} (per-entry evaluation at the common broadcast shape)\\
the computational check & \texttt{polynomialCanonicalEq} $=$ idx-freedom gates $+$ \texttt{Term.toPoly} equality\\
idx-freedom gate (Remark \ref{rem:norm-f7}) & \texttt{Term.idxFree}; \texttt{Term.toPoly}'s wire-keyed variable space\\
typing fields (Remark \ref{rem:norm-b13}) & payload fields \texttt{tyEq}, \texttt{tyDef}, \texttt{finL}, \texttt{finR} (\texttt{PNormData})\\
premise-freedom ($\Phi$ arbitrary) & the \rulename{Norm} convention; \texttt{pre} unused in \texttt{sound\_norm}\\
wire placement & \texttt{RuleTag.norm}\\
\end{anchors}
